\subsection{Ablation Study}
Table~\ref{tab:ablation} reports the complete six-variant ablation results on the CEC2017 benchmark suite. Lower average rank indicates better overall performance. ARPSO-SRR and ARPSO-Local obtain the same best average rank of 2.71, followed by ARPSO-Fixed and ARPSO-EIS. This result indicates that the proposed search-resource reallocation mechanism is competitive with the more aggressive local restart strategy while preserving the simplicity of the adaptive hybrid restart framework.

\begin{table}[!t]
\caption{Six-Variant Ablation Results on the Completed Ablation Functions}
\centering
\scriptsize
\setlength{\tabcolsep}{2.4pt}
\resizebox{\columnwidth}{!}{%
\begin{tabular}{lcccc}
\toprule
Algorithm & Avg. rank & Avg. restarts & Rest. particles & Runtime (s) \\
\midrule
ARPSO-Local  & 2.71 & 3561.7 & 33812.2 & 118.24 \\
ARPSO-SRR    & 2.71 & 2229.8 & 14847.4 & 118.42 \\
ARPSO-EIS    & 3.46 & 2048.1 & 13273.5 & 119.40 \\
ARPSO-Fixed  & 3.46 & 1626.6 & 16265.6 & 116.87 \\
ARPSO-Global & 4.25 & 1538.4 & 10613.0 & 116.63 \\
PSO-RS       & 4.39 & 1324.9 & 8633.6  & 116.13 \\
\bottomrule
\end{tabular}%
}
\label{tab:ablation}
\end{table}


The restart statistics further clarify the difference between ARPSO-SRR and ARPSO-Local. Although both variants reach the best average rank, ARPSO-Local triggers 3561.7 restart events on average and replaces 33812.2 particles, whereas ARPSO-SRR triggers 2229.8 restart events and replaces 14847.4 particles. In other words, ARPSO-Local replaces about 2.28 times as many particles as ARPSO-SRR. This suggests that ARPSO-SRR does not simply rely on a stronger or more aggressive restart operation; instead, it controls the restart intensity and reallocates search resources in a more economical manner.

Fig.~\ref{fig:ablation} visualizes the ablation comparison. The figure supports the same observation as Table~\ref{tab:ablation}: ARPSO-SRR achieves the best overall ranking performance together with ARPSO-Local, while avoiding the substantially larger particle replacement cost of the local restart variant.

\begin{figure}[!t]
    \centering
    \safeincludegraphics[width=\linewidth]{ablation6_average_rank.pdf}
    \caption{Ablation comparison of the six restart variants on the CEC2017 benchmark suite. Lower average rank indicates better performance.}
    \label{fig:ablation}
\end{figure}

The global Friedman test also confirms that the restart design has a statistically observable influence on the overall performance. Across 28 benchmark functions and 6 variants, the Friedman statistic is 21.6596 with $p=6.08\times 10^{-4}$. The pairwise Wilcoxon summary should be interpreted conservatively. ARPSO-SRR obtains 6/22/0 win/no-significant-difference/loss counts against PSO-RS, 8/19/1 against ARPSO-Global, 5/22/1 against ARPSO-EIS, and 1/24/3 against ARPSO-Local. Therefore, the ablation study does not suggest that ARPSO-SRR significantly dominates all restart variants on every function. Instead, it shows that ARPSO-SRR provides a favorable balance between ranking performance and restart cost, which is consistent with the motivation of computational resource reallocation.

